(1) $L_{\rm diff}=d\lambda/dI=L_0/[1+(I/I_s)^2]>0$，所以磁链随电流单调，微分磁能正定；高电流时磁链趋于 $\pm\pi L_0I_s/2$，这只是所给理想化饱和模型的极限。
(2) 场能 $W(I)=\int_0^I i\,d\lambda(i)=L_0I_s^2\ln[1+(I/I_s)^2]/2$。共能 $W_c=\int_0^I\lambda(i)\,di=L_0I_s[I\arctan(I/I_s)-I_s\ln(1+I^2/I_s^2)/2]$。直接相加得到 $W+W_c=I\lambda$。
(3) 法拉第端电压 $v=d\lambda/dt=L_{\rm diff}(I)\dot I$，端口功率 $p=Iv=dW/dt$。在 $|I|\ll I_s$ 时 $\lambda\simeq L_0I-L_0I^3/(3I_s^2)$、$W\simeq L_0I^2/2-L_0I^4/(4I_s^2)$，恢复线性电感。
(4) 无电阻恒压 $V$ 时 $\lambda(t)=Vt$，故 $I(t)=I_s\tan[Vt/(L_0I_s)]$，适用于 $0\le t<\pi L_0I_s/(2V)$（$V>0$）。临界时电流发散说明忽略绕组电阻、饱和外推等假设已失效，不是实际无限能量源。
